\begin{derivation}[四点树级函数的收缩与组合因子]

正文\eqnrefs{eq:phi4-first-dyson-r68} 是 $\varphi^4$ 理论四点函数的一阶 Dyson 插入。连通树级项要求每个外场都与顶点 $z$ 上四个场中的一个且仅一个收缩，因此共有 $4!$ 种相同配对；它们与作用量中的 $1/4!$ 抵消后得到正文\eqnrefs{eq:phi4-tree-green-r3}。



四点函数的一阶 Dyson 项为
\begin{equation*}
 G^{(4,1)}(x_1,x_2,x_3,x_4)
 =-\frac{\ii g_4}{4!\,\hbar^2c^2}
 \int\dd^4 z\,
 \langle0|\Tord\{\varphi_1\varphi_2\varphi_3\varphi_4\varphi^4(z)\}|0\rangle .
\end{equation*}
要得到“一个顶点连接四个外点”的连通树图，每个外部场必须与 $z$ 点四个场中的一个且仅一个收缩。先给 $\varphi_1$ 选 $z$ 点场有 $4$ 种；再给 $\varphi_2$ 选剩余场有 $3$ 种；之后分别有 $2$ 种和 $1$ 种，所以
\begin{equation*}
 N_{\rm contr}=4\times3\times2\times1=4!.
\end{equation*}
每一种收缩都给出同一个核乘积
\begin{equation*}
 \Delta_F(x_1-z)\Delta_F(x_2-z)\Delta_F(x_3-z)\Delta_F(x_4-z).
\end{equation*}
因此 $4!$ 个相同项与相互作用中的 $1/4!$ 抵消，得到
\begin{equation*}
 G_c^{(4)}
 =-\frac{\ii g_4}{\hbar^2c^2}
 \int\dd^4 z\,\prod_{a=1}^4\Delta_F(x_a-z),
\end{equation*}
即正文\eqnrefs{eq:phi4-tree-green-r3}。

再把四条外线分别 Fourier 展开。四个传播子的指数中所有 $z$ 依赖合成
\begin{equation*}
 \exp\!\left[-\frac{\ii}{\hbar}
 (p_1+p_2-p_3-p_4)\cdot z\right].
\end{equation*}
于是顶角位置积分给
\begin{equation*}
 \int\dd^4 z\,
 \ee^{-\ii(p_1+p_2-p_3-p_4)\cdot z/\hbar}
 =(2\pi\hbar)^4\delta^{(4)}(p_1+p_2-p_3-p_4).
\end{equation*}
所以顶角四动量守恒正是局域相互作用对顶角位置积分后的 Fourier $\delta$ 函数。
\end{derivation}

\begin{derivation}[有限时间跃迁概率与黄金法则]
记 $\Delta E=E_f-E_i$。一阶振幅为

\begin{align*}
 c_f^{(1)}(T)
 &=-\frac{\ii}{\hbar}V_{fi}
 \int_{-T/2}^{T/2}\dd t\,\ee^{\ii\Delta E t/\hbar}\\
 &=-\frac{\ii}{\hbar}V_{fi}
 \left[\frac{\hbar}{\ii\Delta E}
 \ee^{\ii\Delta E t/\hbar}\right]_{-T/2}^{T/2}\\
 &=-\frac{2\ii V_{fi}}{\Delta E}
 \sin\!\left(\frac{\Delta E T}{2\hbar}\right).
\end{align*}
把 $\sinc x\equiv\sin x/x$ 提出，得到
\begin{equation*}
 |c_f^{(1)}(T)|^2
 =\frac{|V_{fi}|^2T^2}{\hbar^2}
 \sinc^2\!\left(\frac{\Delta E T}{2\hbar}\right),
\end{equation*}
即\eqnrefs{eq:golden-sinc-r3}。它在有限 $T$ 时不是 $\delta$ 函数；主峰宽度量级是 $\hbar/T$。

对连续末态求和，概率为
\begin{equation*}
 P_{i\to f}(T)
 =\int\dd E\,\nu_f(E)
 \frac{|V_{fi}(E)|^2T^2}{\hbar^2}
 \sinc^2\!\left[\frac{(E-E_i)T}{2\hbar}\right].
\end{equation*}
这里 $\nu_f(E)$ 是末态每单位能量的态密度，SI 量纲为 $\mathrm{J^{-1}}$。有限时间谱核与 Dirac $\delta$ 分布的对应由下列核固定：
\begin{equation*}
 K_T(\Delta E)
 \equiv
 \frac{T}{2\pi\hbar}
 \sinc^2\!\left(\frac{\Delta E\,T}{2\hbar}\right).
\end{equation*}
它具有能量倒数量纲。作变量代换 $x=\Delta E T/(2\hbar)$ 后，
\begin{align*}
 \int_{-\infty}^{+\infty}\dd(\Delta E)\,K_T(\Delta E)
 &=\frac{T}{2\pi\hbar}
 \frac{2\hbar}{T}
 \int_{-\infty}^{+\infty}\dd x\,\sinc^2x\\
 &=\frac1\pi\int_{-\infty}^{+\infty}\dd x\,\sinc^2x=1.
\end{align*}
因此 $K_T$ 是面积固定为一、宽度随 $T^{-1}$ 收缩的归一化谱核；对任意在 $E_i$ 附近平滑的测试函数 $F(E)$，$T\to\infty$ 时满足
\begin{equation*}
 \int\dd E\,F(E)K_T(E-E_i)\longrightarrow F(E_i),
\end{equation*}
也就是 $K_T(\Delta E)\to\delta(\Delta E)$ 的分布极限。
令
\begin{equation*}
F(E)\equiv |V_{fi}(E)|^2\nu_f(E).
\end{equation*}
由 $K_T$ 的定义，有限时间概率可以严格改写为
\begin{equation*}
P_{i\to f}(T)
=\frac{2\pi T}{\hbar}
\int\dd E\,F(E)K_T(E-E_i).
\end{equation*}
再定义有限时间余项
\begin{equation*}
\delta F_T
\equiv
\int\dd E\,[F(E)-F(E_i)]K_T(E-E_i),
\end{equation*}
利用 $\int K_T\dd E=1$ 得到精确等式
\begin{equation*}
P_{i\to f}(T)
=\frac{2\pi T}{\hbar}[F(E_i)+\delta F_T].
\end{equation*}
前面已经证明 $K_T\to\delta$ 的分布极限，所以对满足该近似恒等核收敛条件的连续 $F$，有 $\lim_{T\to\infty}\delta F_T=0$。因此长时间跃迁率严格由极限定义
\begin{equation*}
\lim_{T\to\infty}\frac{P_{i\to f}(T)}{T}
=\frac{2\pi}{\hbar}F(E_i).
\end{equation*}
也就是
\begin{equation*}
 w_{i\to f}=\frac{2\pi}{\hbar}|V_{fi}|^2\nu_f(E_i),
\end{equation*}
即 Fermi 黄金律~\eqnrefs{eq:fermi-golden-r3}。能量 $\delta$ 分布由有限时间谱核与连续谱卷积的长时间极限产生；有限 $T$ 时应保留具有有限宽度的 $\sinc^2$ 核。
\end{derivation}

\begin{derivation}[外腿极点怎样给出散射振幅]
先只处理一条外腿，因为多腿公式只是把同一个极点提取操作重复 $n$ 次。设 $\zeta=p^2-m_{\rm phys}^2c^2$，并把 $\zeta$ 暂时视为复变量。对时序关联函数中与坐标 $x$ 对应的那一个场插入完整态，稳定单粒子态给出
\begin{equation*}
 \langle\Omega|\varphi_H(0)|p,\lambda\rangle
 =Z^{1/2}\,\mathcal N_p,
\end{equation*}
其中 $\mathcal N_p$ 是由全书协变态归一化固定的运动学因子。把该单粒子贡献与其余所有中间态分开，Fourier 变换后必然具有形式
\begin{equation*}
 \widetilde G_{n,c}(p,\ldots)
 =\frac{\ii\hbar^3c\,Z^{1/2}}
 {\zeta}\,
 \widetilde F_{n-1}(p;\ldots)
 +\widetilde R^{(1)}_n(p,\ldots),
\end{equation*}
这里 $\widetilde F_{n-1}$ 是已经把这一条外腿换成归一化单粒子态后的矩阵元；$\widetilde R^{(1)}_n$ 包含连续谱和在这一外部通道中正则的项。稳定粒子极点与多粒子阈值分离，所以在 $\zeta=0$ 的一个小邻域内
\begin{equation*}
 \lim_{\zeta\to0}\zeta\,\widetilde R^{(1)}_n=0.
\end{equation*}
因此极点留数不是一句“外腿传播子可以删掉”，而是如下复分析极限：
\begin{align*}
 \lim_{\zeta\to0}
 Z^{-1/2}\frac{\zeta}{\ii\hbar^3c}
 \widetilde G_{n,c}
 &=\lim_{\zeta\to0}
 \left[
 \widetilde F_{n-1}
 +Z^{-1/2}\frac{\zeta}{\ii\hbar^3c}
 \widetilde R^{(1)}_n
 \right]\\
 &=\widetilde F_{n-1}(p_{\rm on};\ldots).
\end{align*}
Feynman 的 $+\ii0^+$ 是实动量边界值的处方；截肢本身应理解为先在复 $\zeta$ 平面取极点留数，再回到质量壳，因此不会出现把 $\zeta/(\zeta+\ii0^+)$ 当作普通数值比值的歧义。

对第二、第三直到第 $n$ 条外腿重复同一操作。若定义 $\zeta_i=p_i^2-m_i^2c^2$，则联合极点邻域中的多变量 Laurent 展开可写成
\begin{align*}
 \widetilde G_{n,c}
 ={}&(2\pi\hbar)^4\delta^{(4)}\!\left(\sum_i\sigma_i^{\rm io}p_i\right)
 \left[\prod_i
 \frac{\ii\hbar^3c\,Z_i^{1/2}}{\zeta_i}\right]
 \ii\mathfrak A_{fi}(\{p_i\})
 +\widetilde R_n,\\
 &\lim_{\zeta_1,\ldots,\zeta_n\to0}
 \left(\prod_i\zeta_i\right)\widetilde R_n=0.
\end{align*}
第二行正是“余项至少缺少一条完整外腿极点”的数学含义。于是
\begin{align*}
 &\lim_{\zeta_1,\ldots,\zeta_n\to0}
 \left[\prod_iZ_i^{-1/2}
 \frac{\zeta_i}{\ii\hbar^3c}\right]
 \widetilde G_{n,c}\\
 &\qquad=(2\pi\hbar)^4\delta^{(4)}\!\left(\sum_i\sigma_i^{\rm io}p_i\right)
 \ii\mathfrak A_{fi}(\{p_{i,\rm on}\}),
\end{align*}
这就得到正文的动量空间 LSZ 公式。注意截肢和在壳极限的顺序来自极点留数的定义：先乘 $\zeta_i$ 提取留数，再令 $\zeta_i\to0$；若反过来，原 Green 函数正好落在极点上。

现在把这一抽象结构用于树级 $\varphi^4$。一次 Dyson 插入给
\begin{equation*}
 -\frac{\ii}{\hbar c}
 \int\dd^4x\,
 \frac{g_4}{4!\,\hbar c}\,
 \varphi_I(x)^4.
\end{equation*}
把四个外部场分别与顶角处四个场收缩有 $4!$ 种双射，恰好消掉相互作用中的 $1/4!$。每条收缩给一个自由 Feynman 传播子，所以 Fourier 变换后的连通四点函数为
\begin{equation*}
 \widetilde G^{(1)}_{4,c}
 =-\frac{\ii g_4}{\hbar^2c^2}
 (2\pi\hbar)^4\delta^{(4)}(p_1+p_2-p_3-p_4)
 \prod_{a=1}^{4}
 \frac{\ii\hbar^3c}{p_a^2-m^2c^2+\ii0^+}.
\end{equation*}
在树级 $Z=1$。令 $\zeta_a=p_a^2-m^2c^2$，逐腿乘逆极点因子后，乘积中的四条外腿逐一消去：
\begin{align*}
 &\left[\prod_{a=1}^{4}
 \frac{\zeta_a}{\ii\hbar^3c}\right]
 \widetilde G^{(1)}_{4,c}\\
 &\quad=-\frac{\ii g_4}{\hbar^2c^2}
 (2\pi\hbar)^4\delta^{(4)}(p_1+p_2-p_3-p_4)
 \prod_{a=1}^{4}
 \frac{\zeta_a}{\zeta_a+\ii0^+}.
\end{align*}
按上面说明的极点留数顺序取极限，四个因子都贡献单位留数，于是得到正文\eqnrefs{eq:phi4-lsz-result-r7}。因此树级顶角规则不是另记的一条口诀：$4!$ 来自 Wick 收缩计数，四条传播子来自外部收缩，而 LSZ 的四个逆极点恰好把这些外腿传播子截去，只留下局域顶角系数。
\end{derivation}
